\originalchapter{作业六：Maxwell 关系的使用}
\psource{8.044，2013 春，PS6，3 月 20 日截止。中文题面；据官方答案核对的 AI 独立解答。}
\originalsection{1．固体声波的温度扰动}
\begin{exercise}
均匀各向同性固体中小扰动 $\Delta P(x,t)=\Delta P\cos[k(x-vt)]$、$\Delta T(x,t)=\Delta T\cos[k(x-vt)]$；频率高时忽略峰谷间导热，局部过程绝热。原卷省略相位中的波数 $k$，这里补出以使相位无量纲。状态方程 $V=V_0(1+\alpha T-\kappa_TP)$。利用 $U_V|_T=T P_T|_V-P$，求振幅比，表示成 $\alpha,\kappa_T,T,C_V/V_0$ 的组合。
\end{exercise}
\begin{solution}
$P_T|_V=\alpha/\kappa_T$，故 $\delta Q=C_V\dd T+(\alpha T/\kappa_T)\dd V$。代入 $\dd V=V_0(\alpha\dd T-\kappa_T\dd P)$ 并令热零，
\[\frac{\Delta T}{\Delta P}=\frac{\alpha T}{C_V/V_0+\alpha^2T/\kappa_T}.\]
系数在平衡背景处取值；若 $\alpha<0$，压缩温度扰动也可负。公式描述近似流体式体积应变，未包括剪切弹性。
\end{solution}
\originalsection{2．薄膜能量}
\begin{exercise}
膜的广义张力 $\mathcal S(T,A)=-N\kb T/(A-b)$，$A>b$，$b$ 为面积，功 $\mathcal S\dd A$。(a) 证明恒温时恒面积热容 $C_A$ 与面积无关；(b) 证明内能仅依赖温度，并用热容写出。
\end{exercise}
\begin{solution}
$\dd F=-S\dd T+\mathcal S\dd A$ 给 $S_A|_T=-\mathcal S_T=N\kb/(A-b)$，故 $(C_A)_A|_T=T\partial_T S_A=0$。$E_A|_T=\mathcal S-T\mathcal S_T=0$，所以 $E(T)=E(T_*)+\int_{T_*}^T C_A(t)\dd t$。不必假定此模型可靠至零温。
\end{solution}
\originalsection{3．凝聚气体的相容状态方程}
\begin{exercise}
转变以下模型为 $P=aT^{5/2}+bT^3+cV^{-2}$、$C_V=dT^{3/2}V+eT^2V+fT^{1/2}$，系数均恒定。(a) 求 $\dd E(T,V)$；(b) 用全微分条件找系数关系；(c) 求 $E$；(d) 求 $S$。题中以 1995 年实验凝聚为背景，所给模型应作为待检验的现象公式。
\end{exercise}
\begin{solution}
(a) $\dd E=C_V\dd T+[3aT^{5/2}/2+2bT^3-c/V^2]\dd V$。(b) 交叉偏导相等给 $d=15a/4$、$e=6b$，$c,f$ 不受此条件约束。(c)
\[E=\frac32aVT^{5/2}+2bVT^3+\frac23fT^{3/2}+\frac cV+E_* .\]
(d) $\dd S=(C_V/T)\dd T+P_T\dd V$，因此 $S=5aVT^{3/2}/2+3bVT^2+2f\sqrt T+S_*$。常数取决于能量基准与参考熵；相容性本身不证明任意系数满足稳定性。
\end{solution}
\originalsection{4．Curie--Weiss 顺磁体}
\begin{exercise}
$M=AH/(T-T_0)$，仅适用 $T>T_0$。(a) 证明恒温时 $C_M$ 与 $M$ 无关；(b) 以 $C_M(T),A,T_0$ 表示 $E(T,M)$；(c) 求熵。
\end{exercise}
\begin{solution}
由 $H=M(T-T_0)/A$，$S_M|_T=-H_T|_M=-M/A$，所以 $(C_M)_M=T\partial_T S_M=0$。$E_M|_T=H-TH_T=-MT_0/A$，故
\[E=E_*+\int_{T_*}^T C_M(t)\dd t-\frac{T_0M^2}{2A},\quad
S=S_*+\int_{T_*}^T\frac{C_M(t)}t\dd t-\frac{M^2}{2A}.\]
参考 $T_*>T_0$；不能像形式积分那样未经说明从模型不适用的零温起积分。
\end{solution}
